% !TeX root = ../memoire.tex

\chapter*{Remerciements}
\addcontentsline{toc}{chapter}{Remerciements}

Arrivée au terme de ce mémoire, mes premiers remerciements vont à Emmanuelle Bermès, responsable pédagogique du master Technologies numériques appliquées à l'histoire et directrice de ce travail. Je lui suis reconnaissante pour son accompagnement et ses conseils, mais surtout pour la confiance qu'elle m'a accordée et renouvelée au cours de ces dernières années. Au terme d'un parcours sinueux, je mesure aujourd'hui tout ce que cette confiance initiale a rendu possible.

Je remercie très chaleureusement Sophie Dikoff, ma tutrice de stage, et Emma Bertrand, assistante archiviste à la Mission Archives historiques. Je leur suis particulièrement reconnaissante pour la liberté qu'elles m'ont laissée dans mes recherches, leur intérêt constant pour les pistes explorées et la disponibilité avec laquelle elles ont répondu à mes nombreuses questions. Leur connaissance approfondie des fonds et des documents susceptibles de les éclairer a été déterminante. Un travail achevé laisse assez peu voir tout ce qu'il doit aux personnes qui l'ont rendu possible ; celui-ci doit beaucoup à leur expérience, à nos échanges et à leur capacité à m'orienter au sein d'un ensemble documentaire dont je découvrais encore l'étendue.

Ce stage m'a également permis de rencontrer, à l'Université de Montpellier comme dans d'autres établissements, de nombreuses personnes dont les retours d'expérience ont largement nourri ma réflexion. Je remercie tout particulièrement Hélène Lorblanchet, Marie Nikichine, Guenno Marti, Anne-Sophie Gagnal, Audrey Cogoluègnes et Audrey Monet, au SCDI, pour nos échanges autour de Foli@ et, plus largement, pour m'avoir fait découvrir les compétences et les moyens que ce service mobilise en faveur du patrimoine universitaire.

Plusieurs professionnels extérieurs à l'Université ont également accepté de partager leur expérience. Je remercie Pauline Laugraud et Jérémie Léonard, du Musée des Confluences de Lyon, pour leur retour détaillé sur Flora et sur l'articulation entre archives, collections et outils numériques ; Océane Valencia, à Sorbonne Université, ainsi que Sarah Cadorel et Emma Bahous, à l'INALCO, pour leurs éclairages sur les logiciels archivistiques et les projets de réinformatisation ; Céline Dehondt, des Archives départementales de l'Hérault, pour son retour d'expérience sur la structuration et la diffusion des instruments de recherche ; et Claire Garcia, des Archives municipales et métropolitaines de Montpellier, pour la présentation de leur dispositif de diffusion en ligne.

Au sein de l'Université de Montpellier, je remercie également Christophe Salvaing, de la DSIN, et Olivier Hirt, de la Direction de la communication, pour leur disponibilité et les solutions techniques proposées dès les premières réflexions autour du portail ; Pascaline Todeschini, pour la présentation de Calames ; Morgane Villa-Salvignol, du service communication de l'UFR de médecine, pour son intérêt et son enthousiasme à l'égard du projet ; Manon Guillermin, pour son accompagnement lors des expérimentations menées dans Flora ; ainsi que Caroline Ducourau, pour son soutien et ses encouragements. La diversité des métiers, des outils et des expériences rencontrés au cours de ces mois a largement contribué à élargir le projet initial et, avec lui, les perspectives développées dans ce mémoire.

Les quelques lignes qui suivent doivent moins à ce mémoire qu'aux années qui l'ont précédé.

Je remercie Florent Routier qui, depuis bien plus longtemps que moi, semble avoir été certain de ma réussite. Sa confiance ayant déjà fait ses preuves à plus grande échelle, il n'est pas si surprenant qu'elle ait fini par avoir quelque effet sur moi.

Je remercie bien évidemment ma famille, qui a accompagné avec une constance remarquable les différentes étapes de ces dernières années. Sa contribution directe à ce travail aura été plus limitée : je lui ai épargné la maîtrise exacte de l'objet de mes études comme la lecture des pages qui en sont finalement résultées.

Je remercie Nicolas pour aujourd'hui et pour demain.

J'ai une pensée particulière pour ma mère, qui avait assez tôt formé pour moi des ambitions que je n'ai pas toujours accueillies avec le même enthousiasme, mais qui n'aurait certainement pas boudé le plaisir de constater où elles m'ont menée.

Je remercie enfin les membres du jury pour le temps qu'ils consacreront à la lecture de ce travail, dont les proportions trahissent sans doute assez bien mon enthousiasme pour le sujet. J'espère que la richesse des archives montpelliéraines et l'étendue des questions qu'elles soulèvent justifieront le zèle avec lequel j'ai tenté de leur rendre justice. Ce mémoire constitue en tout cas une première étape dans ma découverte de l'histoire de Montpellier, dont ces fonds m'ont surtout fait mesurer l'étendue de ce qu'il reste à explorer.